

\begin{derivation}[从\eqnrefs{eq:hggamma-transverse-tensor-dp11}到\eqnrefs{eq:sm-hgamma-loop-functions-dp41}]
由 Yukawa 相互作用，$h f\bar f$ 的约化顶角为 $-\ii m_fc^2/v_E$，光子顶角为 $-\ii g_{\rm em}Q_f\gamma^\mu$。 两种相反的光子排序对应两张三角图，必须同时保留。 提出共同的外腿归一化因子后，定义约化张量

\begin{align*}
 \mathcal M_f^{\mu\nu}
 ={}&-N_c^f(g_{\rm em}Q_f)^2\frac{m_fc^2}{v_E}
 \int_\ell^{(d)}
 \frac{\Tr\!\left[
 (\slashed\ell+m_fc)\gamma^\mu
 (\slashed\ell+\slashed k_1+m_fc)\gamma^\nu
 (\slashed\ell-\slashed k_2+m_fc)
 \right]}
 {D_0D_1D_2}
 \nonumber\\
 &+(\mu,k_1\leftrightarrow\nu,k_2),
\end{align*}
其中
\begin{align*}
 D_0&=\ell^2-m_f^2c^2+\ii0,\\
 D_1&=(\ell+k_1)^2-m_f^2c^2+\ii0,\\
 D_2&=(\ell-k_2)^2-m_f^2c^2+\ii0,
\end{align*}
两个外部光子满足
\begin{equation*}
 k_1^2=k_2^2=0,
 \qquad
 (k_1+k_2)^2=m_h^2c^2,
 \qquad
 2k_1\!\cdot k_2=m_h^2c^2.
\end{equation*}
闭合费米子环给出该振幅的整体负号。

三个分母用单纯形上的 Feynman 参数合并：
\begin{equation*}
 \frac1{D_0D_1D_2}
 =2\int_0^1\dd x\int_0^{1-x}\dd y\,
 \frac1{[xD_1+yD_2+(1-x-y)D_0]^3}.
\end{equation*}
逐项展开括号：
\begin{align*}
 xD_1+yD_2+(1-x-y)D_0
 ={}&\ell^2+2\ell\!\cdot(xk_1-yk_2)
 -m_f^2c^2
 \\
 &-2xy\,k_1\!\cdot k_2+\ii0.
\end{align*}
令
\begin{equation*}
 L\equiv\ell+xk_1-yk_2,
\end{equation*}
并用 $k_1^2=k_2^2=0$，分母化成
\begin{equation*}
 [L^2-\Delta_f(x,y)+\ii0]^3,
 \qquad
 {\Delta_f=m_f^2c^2-xy\,m_h^2c^2.}
\end{equation*}

接着把分子真正展开。为避免在长迹中反复书写，记
\[
a\equiv\ell,\qquad b\equiv\ell+k_1,\qquad c\equiv\ell-k_2,
\qquad M_f\equiv m_fc.
\]
第一种光子排序的迹为
\begin{align*}
N^{\mu\nu}
&\equiv\Tr\!\left[(\slashed a+M_f)\gamma^\mu
(\slashed b+M_f)\gamma^\nu(\slashed c+M_f)\right].
\end{align*}
含零个或两个 $M_f$ 的项分别含奇数个 $\gamma$ 矩阵，迹为零；只剩一个 $M_f$ 和三个 $M_f$ 的项。逐项使用
\[
\Tr(\gamma^\alpha\gamma^\beta)=4\eta^{\alpha\beta},\qquad
\Tr(\gamma^\alpha\gamma^\beta\gamma^\gamma\gamma^\delta)
=4(\eta^{\alpha\beta}\eta^{\gamma\delta}
-\eta^{\alpha\gamma}\eta^{\beta\delta}
+\eta^{\alpha\delta}\eta^{\beta\gamma}),
\]
得到
\begin{align*}
\frac{N^{\mu\nu}}{4M_f}
={}&b^\mu c^\nu+c^\mu b^\nu-\eta^{\mu\nu}b\!\cdot c
+a^\mu c^\nu-a^\nu c^\mu+\eta^{\mu\nu}a\!\cdot c\\
&+a^\mu b^\nu+a^\nu b^\mu-\eta^{\mu\nu}a\!\cdot b
+M_f^2\eta^{\mu\nu}.
\end{align*}
第二张三角图由 $(\mu,k_1)\leftrightarrow(\nu,k_2)$ 得到。于是两张图相加以后，整个问题已经从六个 $\gamma$ 矩阵的迹化成普通 Lorentz 张量代数。

平移到 $L=\ell+xk_1-yk_2$ 后，奇数次 $L$ 的积分为零；秩二项用
\begin{equation*}
\int\dd^dL\,L^\mu L^\nu F(L^2)
=\frac{\eta^{\mu\nu}}{d}\int\dd^dL\,L^2F(L^2)
\end{equation*}
约化。为了不在中间步骤预先假定横向性，定义正文张量
\[
T^{\mu\nu}\equiv k_1\!\cdot k_2\,\eta^{\mu\nu}-k_2^\mu k_1^\nu,
\]
其自收缩满足
\begin{equation*}
T_{\mu\nu}T^{\mu\nu}=(d-2)(k_1\!\cdot k_2)^2.
\end{equation*}
因此横向系数可以在所有张量积分完成后由
\begin{equation*}
\mathcal C_f
=\frac{T_{\mu\nu}\mathcal M_f^{\mu\nu}}
{(d-2)(k_1\!\cdot k_2)^2}
\end{equation*}
直接投影。把上面的显式 $N^{\mu\nu}$ 代入这一投影，所有 $L^2$ 项都通过
\begin{equation*}
\frac{L^2}{(L^2-\Delta_f)^3}
=\frac1{(L^2-\Delta_f)^2}
+\frac{\Delta_f}{(L^2-\Delta_f)^3}
\end{equation*}
化成二点型与三点型标量积分；两种光子排序相加后，二点型局域项彼此抵消。剩余三点项为
\begin{equation*}
 A_{1/2}(\tau)
 =4\tau\int_0^1\dd x\int_0^{1-x}\dd y\,
 \frac{1-4xy}{\tau-4xy-\ii0},
 \qquad
 \tau\equiv\frac{4m_f^2}{m_h^2}.
\end{equation*}
这里的分子 $1-4xy$ 是费米子迹中质量项与动量项共同作用的结果；只保留其中任一部分都会破坏重质量极限。

双积分可用下列组合化为闭式：
\begin{equation*}
 I(\tau)
 \equiv\int_0^1\dd x\int_0^{1-x}\dd y\,
 \frac1{\tau-4xy-\ii0}.
\end{equation*}
先对 $y$ 积分：
\begin{align*}
 I(\tau)
 &=\int_0^1\frac{\dd x}{4x}
 \ln\!\left[
 \frac{\tau}{\tau-4x(1-x)-\ii0}
 \right].
\end{align*}
再令 $x=(1-z)/2$，使 $4x(1-x)=1-z^2$。对 $\tau\ge1$，积分路径不穿过分支割线；完成变量代换并积分后得到
\begin{equation*}
 I(\tau)=\frac12\arcsin^2\!\frac1{\sqrt\tau}.
\end{equation*}
对 $\tau<1$，分母在积分区域内穿过双费米子阈值；按 $-\ii0$ 处方解析延拓后出现对数与虚部。因此统一定义
\begin{equation*}
 {
 f(\tau)=2I(\tau)=
 \begin{cases}
 \arcsin^2(1/\sqrt\tau),&\tau\ge1,\\[4pt]
 -\dfrac14\left[
 \ln\dfrac{1+\sqrt{1-\tau}}{1-\sqrt{1-\tau}}-\ii\pi
 \right]^2,&\tau<1.
 \end{cases}}
\end{equation*}
$\tau=1$ 正是 $m_h=2m_f$ 的粒子对产生分支点。

现在把上式的分子写成
\begin{equation*}
 1-4xy=(\tau-4xy)+(1-\tau).
\end{equation*}
于是
\begin{align*}
 A_{1/2}(\tau)
 &=4\tau\left[
 \int_\triangle\dd x\dd y
 +(1-\tau)I(\tau)
 \right],\\
 &=4\tau\left[\frac12+\frac{1-\tau}{2}f(\tau)\right],\\
 &={2\tau[1+(1-\tau)f(\tau)]}.
\end{align*}
因此 $A_{1/2}(\tau)$ 恰好化为\eqnrefs{eq:sm-hgamma-loop-functions-dp41} 的第一行。

重费米子极限也可以直接从参数积分检验。$\tau\gg1$ 时
\begin{equation*}
 f(\tau)=\frac1\tau+\frac1{3\tau^2}+\order(\tau^{-3}),
\end{equation*}
代回得到
\begin{equation*}
 {A_{1/2}(\tau)=\frac43+\order(\tau^{-1}).}
\end{equation*}
\end{derivation}

\begin{derivation}[从\eqnrefs{eq:hggamma-transverse-tensor-dp11}到\eqnrefs{eq:sm-hgamma-loop-functions-dp41}]
$W$ 环的困难不在最后一个参数积分，而在自旋一传播子的纵向分子会把圈动量幂次显著抬高。采用幺正规范后，Goldstone 与 ghost 不再作为独立内部线出现，而物理在壳振幅保持规范无关。记
\[
M_W^2\equiv m_W^2c^2,\qquad M_h^2\equiv m_h^2c^2,
\]
则 $W$ 传播子与 $hW^+W^-$ 约化顶角分别为
\begin{equation*}
 D_W^{\alpha\beta}(q)
 =\frac{-\ii}{q^2-M_W^2+\ii0}
 \left(\eta^{\alpha\beta}-\frac{q^\alpha q^\beta}{M_W^2}\right),
 \qquad
 V_{hWW}^{\alpha\beta}=\ii\frac{2(m_Wc^2)^2}{v_E}\eta^{\alpha\beta}.
\end{equation*}
两个 $WW\gamma$ 顶角由正文已建立的 Yang--Mills 顶角给出，另有一个 $WW\gamma\gamma$ 接触顶角。因此幺正规范下需要三角图、交换两个光子的交叉三角图以及接触图，缺少其中任一项都会破坏光子 Ward 恒等式。

令两个光子均出射，$k_1^2=k_2^2=0$，$2k_1\!\cdot k_2=M_h^2$；取三角图的三个分母
\begin{equation*}
D_0=p^2-M_W^2+\ii0,\qquad
D_1=(p-k_1)^2-M_W^2+\ii0,\qquad
D_2=(p-k_1-k_2)^2-M_W^2+\ii0.
\end{equation*}
提出共同因子
\begin{equation*}
\mathcal N_W\equiv\frac{g_{\rm em}^2g}{16\pi^2m_Wc},
\end{equation*}
并暂时不收缩外部偏振。把三张图的 Lorentz 指标全部收缩后，所有含 $k_1^\mu$ 或 $k_2^\nu$ 的项在物理外腿上消失；其余部分可以在五个张量上展开为
\begin{equation*}
\mathcal M_W^{\mu\nu}
=\int_p^{(d)}\!
\left[
\mathscr M_1\eta^{\mu\nu}
+\mathscr M_2p^\mu p^\nu
+\mathscr M_3p^\mu k_1^\nu
+\mathscr M_4k_2^\mu p^\nu
+\mathscr M_5k_2^\mu k_1^\nu
\right].
\end{equation*}
五个系数先分别完成圈积分与张量约化；把同一规范不变图集合求和后，它们自动重组为满足 Ward 恒等式的横向张量。

按正文\eqnrefs{eq:pv-scalar-masters-dp90,eq:pv-rank2-decomposition-dp90} 做 Passarino--Veltman 约化。为缩短公式，以下记
\begin{equation*}
 A_0\equiv A_0(M_W^2),\qquad
 B_h\equiv B_0(M_h^2;M_W^2,M_W^2),\qquad
 B_0\equiv B_0(0;M_W^2,M_W^2),
\end{equation*}
\begin{equation*}
 C_h\equiv C_0(M_h^2,0,0;M_W^2,M_W^2,M_W^2).
\end{equation*}
五类张量积分分别化为
\begin{align*}
\int_p^{(d)}\mathscr M_1\eta^{\mu\nu}
={}&\mathcal N_W
\Bigl[4M_W^2(1-2M_h^2C_h)-(M_h^2+6M_W^2)B_h\Bigr]\eta^{\mu\nu},\\[2mm]
\int_p^{(d)}\mathscr M_2p^\mu p^\nu
={}&\frac{\mathcal N_W}{M_h^2}
\Bigl\{
(M_h^2+6M_W^2)(1+B_h+2M_W^2C_h)
(M_h^2\eta^{\mu\nu}-2k_2^\mu k_1^\nu)\\
&\hspace{18mm}-4M_W^2M_h^2\eta^{\mu\nu}
+2(M_h^2+6M_W^2)(2B_0-B_h)k_2^\mu k_1^\nu
\Bigr\},\\[2mm]
\int_p^{(d)}\mathscr M_3p^\mu k_1^\nu
={}&\frac{\mathcal N_W}{4M_W^2M_h^2}
\Bigl[
16M_W^2(M_h^2+6M_W^2)
-2(7M_h^2+48M_W^2)A_0\\
&\hspace{22mm}+(96M_W^4+2M_W^2M_h^2-M_h^4)B_h
\Bigr]k_2^\mu k_1^\nu,\\[2mm]
\int_p^{(d)}\mathscr M_4k_2^\mu p^\nu
={}&\frac{\mathcal N_W}{4M_W^2}
\Bigl[16M_W^2-18A_0+(M_h^2+14M_W^2)B_h\Bigr]k_2^\mu k_1^\nu,\\[2mm]
\int_p^{(d)}\mathscr M_5k_2^\mu k_1^\nu
={}&\mathcal N_W\Bigl[4B_0+16M_W^2C_h\Bigr]k_2^\mu k_1^\nu.
\end{align*}
这些式子展示了为什么不能在积分前把 $p^\mu p^\nu$ 直接替换为 $\eta^{\mu\nu}p^2/4$：各项分别含紫外发散的 $A_0$ 与 $B_0$，只有维数正规化下完成同一张量约化后，有限常数项才有唯一意义。

现在真正执行求和。相同质量的零动量二点函数满足
\begin{equation*}
 B_0(0;M_W^2,M_W^2)=\frac{A_0(M_W^2)}{M_W^2}-1.
\end{equation*}
把它代入上面五项，先收集 $\eta^{\mu\nu}$ 系数，再收集 $k_2^\mu k_1^\nu$ 系数。$k_2^\mu k_1^\nu$ 前的 $A_0$ 系数为
\begin{align*}
&\frac{4(M_h^2+6M_W^2)}{M_h^2M_W^2}
-\frac{7M_h^2+48M_W^2}{2M_h^2M_W^2}
-\frac{9}{2M_W^2}+\frac{4}{M_W^2}\\
&\qquad
=\frac{8(M_h^2+6M_W^2)-(7M_h^2+48M_W^2)-9M_h^2+8M_h^2}
{2M_h^2M_W^2}=0,
\end{align*}
而 $B_h$ 系数同样逐项给出
\begin{align*}
&-\frac{4(M_h^2+6M_W^2)}{M_h^2}
-\frac{M_h^4-2M_h^2M_W^2-96M_W^4}{4M_h^2M_W^2}
+\frac{M_h^2+14M_W^2}{4M_W^2}\\
&\qquad
=\frac{-16M_W^2(M_h^2+6M_W^2)
-(M_h^4-2M_h^2M_W^2-96M_W^4)
+M_h^2(M_h^2+14M_W^2)}{4M_h^2M_W^2}=0.
\end{align*}
因此所有一、二点母积分和紫外局域项都从物理振幅中消失。剩下的有限张量恰好组合为
\begin{equation*}
\mathcal M_W^{\mu\nu}
=\frac{\mathcal N_W}{M_h^2}
\Bigl[M_h^2+6M_W^2-6M_W^2(M_h^2-2M_W^2)C_h\Bigr]
\left(M_h^2\eta^{\mu\nu}-2k_2^\mu k_1^\nu\right).
\end{equation*}
由于
\begin{equation*}
M_h^2\eta^{\mu\nu}-2k_2^\mu k_1^\nu
=2\left(k_1\!\cdot k_2\,\eta^{\mu\nu}-k_2^\mu k_1^\nu\right),
\end{equation*}
两个光子 Ward 恒等式已经显式满足；横向性是全部图求和后的结果，而不是在单图上强加的投影。

最后只剩三点标量母积分。对三个等质量分母作 Feynman 参数化，正文\eqnrefs{eq:pv-scalar-masters-dp90} 给出
\begin{align*}
 C_h
&=-\int_0^1\dd x\int_0^{1-x}\dd y\,
\frac{1}{M_W^2-xyM_h^2-\ii0}\\
&=-\frac{2}{M_h^2}f(\tau_W),
\qquad
\tau_W\equiv\frac{4M_W^2}{M_h^2}=\frac{4m_W^2}{m_h^2}.
\end{align*}
代回有限系数并除去公共横向张量，得到
\begin{align*}
\frac{2}{M_h^2}
\Bigl[M_h^2+6M_W^2-6M_W^2(M_h^2-2M_W^2)C_h\Bigr]
&=2+3\tau_W+3\tau_W(2-\tau_W)f(\tau_W).
\end{align*}
按照正文对一圈函数的振幅符号约定，因此
\begin{equation*}
\boxed{
A_1(\tau)
=-\left[2+3\tau+3\tau(2-\tau)f(\tau)\right],}
\end{equation*}
即\eqnrefs{eq:sm-hgamma-loop-functions-dp41} 的第二行。用
$f(\tau)=\tau^{-1}+(3\tau^2)^{-1}+\order(\tau^{-3})$
还可立即检验
\begin{equation*}
A_1(\tau)=-7+\order(\tau^{-1}).
\end{equation*}
\end{derivation}
